\documentclass[../../main.tex]{subfiles}
\begin{document}


{\bf [解]}
(1)
まず二変数の相加相乗平均不等式（AM-GM）を示す．
\begin{align}
 \dfrac{a+b}{2} \geqq \sqrt{ab} \label{eq:1}
\end{align}
両辺正だから二乗してよく
\begin{align}
 \text{\cref{eq:1}} \iff (a+b)^2 \geqq 4ab \iff (a-b)^2 \geqq 0
\end{align}
である．よって\cref{eq:1}は成立し，等号成立は $a=b$ の時である．


次に題意の二つ目の不等式，$N=4$のAM-GMを示す．
$(a,b) = (a+b, c+d)$として\cref{eq:1}に代入すると
\begin{align}
 \dfrac{a+b+c+d}{2} 
     & \geqq \sqrt{(a+b)(c+d)} \\
     & \geqq \sqrt{2\sqrt{ab}\cdot 2\sqrt{cd}} \quad (\because \text{\cref{1}}) \\
     & = 2 \sqrt[4]{abcd} \\
\therefore 
  \dfrac{a+b+c+d}{4} & \geqq \sqrt[4]{abcd} \label{eq:2}
\end{align}
となり示された．等号成立は$a+b=c+d, a=b, c=d$すなわち$a=b=c=d$の時である．

最後に題意の一つ目の不等式，$N=3$のAM-GMを示す．
\cref{eq:2}に$d = \dfrac{a+b+c}{3}$ を代入すると
\begin{align}
    \dfrac{a+b+c+\dfrac{a+b+c}{3}}{4} & \geqq \sqrt[4]{abc \cdot \dfrac{a+b+c}{3}} \\
\therefore 
    \dfrac{a+b+c}{3} & \geqq \sqrt[4]{abc \cdot \dfrac{a+b+c}{3}} 
\end{align}
両辺正だから4乗して整理すると
\begin{align}
    \left(\dfrac{a+b+c}{3}\right)^3 & \geqq abc \\ 
\therefore 
    \dfrac{a+b+c}{3} & \geqq \sqrt[3]{abc} \label{eq:3}
\end{align}
となり示された．等号成立は$a=b=c$の時である．

(2)
まず$P, R$の大小を比較する．$P$に\cref{eq:1}を用いて
\begin{align}
 P 
   &= \dfrac{3a+3b+3c+3d}{12} \\
   &= \dfrac{(a+b)+(a+c)+(a+d)+(b+c)+(b+d)+(c+d)}{12} \\
   &\geqq \dfrac{\sqrt{ab} + \sqrt{ac} + \sqrt{ad} + \sqrt{bc} + \sqrt{bd} + \sqrt{cd}}{6} \\
   &= R \\
\therefore
  P & \geqq R \label{eq:4}
\end{align}
である．

次に$R,S$の大小を比較する．$R$に\cref{eq:3}を用いて
\begin{align}
  R 
    &= \dfrac{2\sqrt{ab} + 2\sqrt{bc} + 2\sqrt{ac} + 2\sqrt{ad} + 2\sqrt{bd} + 2\sqrt{cd}}{12} \\
    &= \dfrac{\left(\sqrt{ab}+\sqrt{bc}+\sqrt{ca}\right)+\left(\sqrt{ab}+\sqrt{bd}+\sqrt{da}\right)+\left(\sqrt{ac}+\sqrt{cd}+\sqrt{da}\right)+\left(\sqrt{bc}+\sqrt{cd}+\sqrt{db}\right)}{12} \\
    & \geqq \dfrac{\sqrt[3]{abc} + \sqrt[3]{abd} + \sqrt[3]{acd} + \sqrt[3]{bcd}}{3} \\
    &= S \\
\therefore
  R & \geqq S \label{eq:5}
\end{align}
である．

最後に$S, Q$の大小を比較する．\cref{eq:2}より
\begin{align}
  S \geqq \sqrt[4]{abcd} = Q \label{eq:6}
\end{align}
である．

以上\cref{eq:4,eq:5,eq:6}から
\begin{align}
    P \geqq R \geqq S \geqq Q 
\end{align}
である．

\end{document}
